\chapter{Introduction}

\section{Contexte}
Dans le cadre de la transformation numérique des établissements d'enseignement supérieur, l'université souhaite
disposer d'une plateforme gérant l'ensemble du cycle de vie d'un stage : publication des offres par les entreprises
partenaires, candidatures des étudiants, sélection, puis suivi du stage jusqu'à son évaluation finale.

Ces activités mobilisent des acteurs et des référentiels distincts (scolarité, relations entreprises, bureau des
stages) qui évoluent à des rythmes différents. Une \textbf{architecture orientée services} (SOA) est donc pertinente :
chaque domaine métier est exposé sous la forme d'un service autonome, au contrat explicite, et un
\textbf{processus métier BPMN} orchestre ces services pour automatiser le traitement d'une candidature, depuis son
dépôt jusqu'à la décision finale.

\section{Objectifs}
Le projet répond aux exigences suivantes du cahier des charges :
\begin{itemize}
  \item trois services \textbf{SOAP} décrits en WSDL : gestion des étudiants, des entreprises et des offres de stage ;
  \item trois services \textbf{REST} : gestion des candidatures, notifications, suivi des stages ;
  \item un \textbf{processus BPMN} d'orchestration comportant tâches de service, événements de début et de fin,
        passerelles de décision et gestion des erreurs ou indisponibilités de services ;
  \item la validation des données, la gestion des erreurs et la protection des échanges ;
  \item en bonus, une \textbf{interface graphique} pour les étudiants, les entreprises et les responsables.
\end{itemize}

\section{Synthèse de la réalisation}
\begin{table}[htbp]
\centering
\caption{Couverture du cahier des charges}
\begin{tabularx}{\textwidth}{@{}lYl@{}}
\toprule
\textbf{Exigence} & \textbf{Réalisation} & \textbf{Chapitre} \\
\midrule
3 services SOAP + WSDL & JAX-WS (Metro), 24 opérations, WSDL publiés sous \code{/ws/*?wsdl} & \ref{chap:soap} \\
3 services REST & JAX-RS (Jersey), 23 points d'accès JSON (+ 7 pour l'orchestration), codes HTTP normalisés & \ref{chap:rest} \\
Persistance & JDBC sur H2 (défaut) ou PostgreSQL, testé sur les deux & \ref{chap:donnees} \\
Processus BPMN & BPMN 2.0 exécuté par Flowable, 4 couloirs, 3 sous-processus, 6 passerelles & \ref{chap:bpmn} \\
Erreurs et indisponibilités & Reprise exponentielle, événements d'erreur, tâche d'incident, simulateur de pannes & \ref{sec:erreurs} \\
Sécurité & WS-Security (SOAP), clé d'API ou HTTP Basic (REST), PBKDF2, masquage des données personnelles, contrôle par rôle & \ref{chap:securite} \\
Notifications & Vrais e-mails HTML envoyés en SMTP et capturés par la boîte de test Mailpit & \ref{chap:rest} \\
Règles de gestion & 32 règles formalisées et implémentées & \ref{chap:regles} \\
Tests & 28 tests d'intégration JUnit, collection Postman (45 assertions), requêtes SoapUI & \ref{chap:tests} \\
Bonus Front-End & Application web à trois espaces aux couleurs de l'INPT, visualisation BPMN en direct & \ref{chap:interface} \\
\bottomrule
\end{tabularx}
\end{table}

\section{Organisation du rapport}
Le chapitre~\ref{chap:architecture} présente l'architecture SOA et les interactions entre services ; le
chapitre~\ref{chap:donnees} le modèle de données. Les chapitres~\ref{chap:soap} et~\ref{chap:rest} documentent les
opérations SOAP et les API REST, le chapitre~\ref{chap:securite} la sécurité. Le chapitre~\ref{chap:bpmn} détaille le
processus BPMN et la gestion des erreurs, le chapitre~\ref{chap:regles} les règles de gestion, le
chapitre~\ref{chap:tests} les scénarios de test et le chapitre~\ref{chap:interface} l'interface graphique.
